\documentclass[12pt,a4paper]{article}
\usepackage[utf8]{inputenc}
\usepackage[T1]{fontenc}
\usepackage{amsmath,amssymb}
\usepackage[margin=2.5cm]{geometry}
\usepackage{setspace}
\usepackage{hyperref}

\onehalfspacing

\begin{document}

\begin{center}
    {\LARGE\bfseries Thermal Spectroscopy of the Minesweeper Probe:\\[4pt]
    Mapping the Phase Transition in LLM Combinatorial Accuracy\\[4pt]}

    \vspace{0.5em}
    {\large Franklin Silveira Baldo}\\
    \vspace{0.5em}
    March 2026
\end{center}

\begin{abstract}
\noindent
[STUB — to be written when experiment is executed.]
This paper maps how the Rosencrantz protocol's measurement accuracy degrades with sampling temperature. Does the model's combinatorial accuracy on Minesweeper boards degrade gracefully or catastrophically? Is there a sharp phase transition? The thermal robustness hierarchy predicts that structural relationships (relative mine probabilities between cells) survive temperatures that destroy absolute accuracy. The experiment runs the three-universe protocol at $T \in \{0.0, 0.1, 0.3, 0.5, 0.7, 1.0, 1.3, 1.5, 2.0\}$ and maps the full landscape.
\end{abstract}

% Research agenda item: Thermal Spectroscopy
% See EXPERIENCE.md "Current Research Agenda" → Near-term
% Protocol: existing minesweeper_basic.py with temperature sweep

\end{document}
